\documentclass[11pt]{article}
\usepackage[margin=1in]{geometry}
\usepackage{hyperref}

\title{Implementation Statement}
\author{Nguyen Huy An Khanh}
\date{COMP2050 Programming Project\\Instructor: Leandro S. Marcolino\\June 21, 2026}

\begin{document}

\maketitle

\section*{Project Overview}

This statement describes how the implementation work for the Minecraft sword PvP agent was performed. The submitted code implements a Monte Carlo Tree Search based combat agent, a Mineflayer server-side bot, a local prismarine-physics duel simulator, experiment logging, and report-material generation.

\section*{Existing Libraries and External Sources}

The project uses open-source npm libraries as infrastructure:
\begin{itemize}
  \item \texttt{mineflayer} is used to connect a bot to a Minecraft server, read nearby player entities, equip a sword, control movement, and perform attacks.
  \item \texttt{prismarine-physics} is used to simulate Minecraft-like player movement in the local duel environment.
  \item \texttt{minecraft-data} and \texttt{prismarine-block} are used to construct the flat local world needed by the physics simulator.
  \item \texttt{vec3} is used for three-dimensional position and velocity vectors.
  \item TypeScript and \texttt{tsx} are used for development and execution.
\end{itemize}

These libraries were not written by me. They provide the environment, Minecraft protocol interaction, block data, physics approximation, and vector primitives. The project does not copy a complete MCTS implementation from an external tutorial or repository.

\section*{AI-Assisted Work}

AI assistance was used during development to help inspect the assignment brief, review and debug TypeScript code, improve the structure of the experiment generator, generate LaTeX report text, and improve readability. AI assistance also helped document and refine the later experimental variations: hybrid MCTS with a heuristic action bonus, risk-aware reward shaping, high-exploration MCTS, and higher-budget MCTS.

The AI-assisted work should be understood as programming and writing assistance, not as an external source of experimental results. The reported results were generated by running the local code in this repository. The manual implementation work described below was written and integrated by me.
\section*{My Own Implementation Work}

My implementation work includes:
\begin{itemize}
\item Designing the Minecraft sword-duel problem and choosing MCTS as the main AI algorithm.
\item Manually writing the Mineflayer-side runtime in \texttt{src/bot.ts}, including server connection configuration, target selection, sword equipping, movement controls, and attack calls.
\item Defining the combat action space and implementing the simulated fighter state, local duel environment, attack cooldowns, health tracking, and duel loop.
\item Implementing the MCTS search process over combat actions, including UCT selection, rollout simulation, backpropagation, and action selection.
\item Designing the reward function and algorithm variations, and comparing them experimentally.
\item Creating the experiment logging pipeline, interpreting the results, and writing the final report.
\end{itemize}

\section*{How to Identify the Main Code}

The main source files are:
\begin{itemize}
  \item \texttt{src/bot.ts}: runtime entry point, Mineflayer server mode, command-line configuration, and control application.
  \item \texttt{src/combat.ts}: combat state, local physics environment, MCTS planner, reward function, algorithm variations, and duel recording.
  \item \texttt{src/report-materials.ts}: repeated experiment generator, CSV summaries, raw JSON traces, and visual material generation.
\end{itemize}

\section*{Academic Integrity Declaration}

To the best of my knowledge, this statement accurately describes the origin of the submitted work. Existing libraries are acknowledged above, AI-assisted work is disclosed, and the project-specific design, integration, experimentation, and interpretation are presented as my own work.

\end{document}
